\section{Uniform truth and finite auxiliary presentations}
\label{sec:p03-admissible-satisfaction}

A truth definition for a structure that is a set is different from a
truth predicate for the entire universe of set theory. This distinction
allows infinitary satisfaction to be computed inside an admissible set,
and allows a recursively axiomatized auxiliary structure to be described
by a single sentence in a logic that can specify the standard natural
numbers. Neither construction asserts that the universe contains its own
full truth predicate.

\begin{hbtheorem}{thm:p03-adm-satisfaction}{proved=annex:p3.adm.12}{Uniform infinitary satisfaction}
Over \(\mathrm{KPU}\) with \(\omega\) as a set, satisfaction for a
set-sized relational structure and a legal coded infinitary formula is
\(\Delta_1\), uniformly in the structure, formula, and assignment.
Formula codes have finitely many free variables and set-sized immediate
constituent families, as in \cref{def:p03-adm-setting}. The evaluation
exists internally and is unique. In particular, if a transitive
admissible set contains the structure and formula code, it computes
their actual satisfaction relation correctly. Empty structures and
nullary predicates are included.
\end{hbtheorem}

\begin{hbdefinition}{def:p03-standard-number-presentation}{A standard-number presentation}
A logic admits a standard-number presentation if a sentence in a finite
expanded signature specifies a distinguished domain \(N\), zero, and
successor, whose induced structure is exactly \((\omega,0,n\mapsto
n+1)\), up to isomorphism, and has such an expansion. Auxiliary objects
outside \(N\) are permitted. Ordinary first-order constraints can be
conjoined with this sentence, and it can be relativized to a named domain.
The quantifier \(Q_0x\,\phi\) means that infinitely many objects satisfy
\(\phi\).
\end{hbdefinition}

\begin{hblemma}{thm:p03-standard-numbers}{proved=annex:p3.adm.13}{Specifying the natural numbers}
Both first-order logic with \(Q_0\) and weak second-order logic admit a
standard-number presentation. The latter quantifies over finite subsets
of its domain. The presentation uses a nonempty linear order without a
largest element in which every proper initial segment is finite.
\end{hblemma}

\begin{hbtheorem}{thm:p03-recursive-auxiliary-presentation}{proved=annex:p3.adm.14}{A finite presentation of a recursive theory}
Let \(T\) be a recursively axiomatized first-order theory in a finite
relational signature, with nonempty structures. In any logic admitting a
standard-number presentation there is a single sentence in a finite
expanded signature whose distinguished reducts are exactly the models
of \(T\). The witnessing expansion may have additional domains.
The construction is uniform in a recursive enumeration of the axioms.
In particular it applies to \(\mathrm{KPU}\) with Infinity and a set of
all urelements, after relationalizing its finite signature. A further
finite constraint restricts its reducts exactly to models whose internal
\(\omega\) is standard.

The additional truth relation is external to the distinguished
\(T\)-structure and is not among the symbols whose truth it evaluates.
No element of that structure is required to code its full satisfaction
relation. This finite auxiliary presentation therefore does not assert
a finite axiomatization of \(T\) in its original language.
\end{hbtheorem}
